\documentclass[12pt]{article}
\usepackage[utf8]{inputenc}
\usepackage[T1]{fontenc}
\usepackage[margin=1in]{geometry}
\usepackage{setspace}
\usepackage{xcolor}
\usepackage[colorlinks=true, urlcolor=blue!60!black, linkcolor=blue!60!black]{hyperref}
\setlength{\parindent}{2em}
\setlength{\parskip}{1.25em}
\renewcommand{\baselinestretch}{1.0}
\begin{document}
\begin{center}
{\LARGE \textit{Hein v. Freedom From Religion Foundation, Inc.}, 551 U.S. 587 (2007)\par}
\vspace{0.8em}
{\large Case Brief --- Four-Page Edition\par}
\end{center}
\vspace{1em}

\section{Memory Jogger}

Federal taxpayers challenging White House faith-based conferences lacked standing, 5--4, because \textit{Flast v.\ Cohen}'s exception covers specific congressional exercises of the Taxing and Spending Clause, not discretionary Executive Branch spending from lump-sum appropriations. Justice Alito's plurality left \textit{Flast} as it found it; Justice Scalia would have overruled it (\href{https://supreme.justia.com/cases/federal/us/551/587/}{551 U.S. 587}).

\section{Detailed Case Facts}

In 2001 President George W. Bush created the White House Office of Faith-Based and Community Initiatives by executive order to help religious and other community groups compete on a level playing field for federal assistance; other orders created centers in several agencies. Congress enacted no statute specifically authorizing or funding these offices; they were paid for from general, lump-sum executive appropriations.

The Freedom From Religion Foundation and three members sued the offices' directors, including Jay Hein, alleging that conferences at which the President and other officials praised faith-based programs and singled them out as worthy of funding told believers they were ``insiders'' and nonbelievers ``outsiders.'' Their only asserted basis for standing was their status as federal taxpayers.

\section{Procedural History}

The Western District of Wisconsin dismissed for lack of standing, confining \textit{Flast} to congressional spending. A divided Seventh Circuit panel, per Judge Posner, reversed, allowing suit whenever the marginal cost of a violation exceeds zero; Judge Ripple dissented. Rehearing en banc was denied 7--4, with Chief Judge Flaum and Judge Easterbrook inviting Supreme Court clarification. The Court reversed on June 25, 2007.

\section{Judicial Votes}

Reversed, \textbf{5--4}, with no majority opinion.
\begin{itemize}
\item \textbf{Plurality:} Justice Alito, joined by Chief Justice Roberts and Justice Kennedy.
\item \textbf{Concurring:} Justice Kennedy.
\item \textbf{Concurring in the judgment:} Justice Scalia, joined by Justice Thomas.
\item \textbf{Dissenting:} Justice Souter, joined by Justices Stevens, Ginsburg, and Breyer.
\end{itemize}

\section{Holding}

Under \textit{Marks}, the plurality controls: \textit{Flast}'s Establishment Clause exception does not extend to purely discretionary executive expenditures funded by general appropriations. A taxpayer must challenge spending expressly authorized or mandated by a specific congressional enactment under Article~I, \S~8. Suggestions that other plaintiffs could challenge extreme abuses, Footnote~7's treatment of committee-report earmarks, and Kennedy's and Scalia's broader views are dicta.

\section{Analysis of Opinions}

\textbf{Plurality (Alito).} Under \textit{Frothingham} and \textit{Doremus}, a taxpayer's interest in the Treasury is shared, minute, and indeterminable. \textit{Flast}'s two-part nexus test has been applied with ``rigor'' and confined to its facts: there, Congress expressly appropriated funds for religious schools, and in \textit{Bowen v.\ Kendrick} the statute contemplated religious grantees. Extending \textit{Flast} to executive spending would reach virtually every executive act and make courts ``continuing monitors'' of the President. The plurality rejected Posner's zero-cost test and declined to reconsider \textit{Flast}.

\textbf{Kennedy (concurring).} \textit{Flast} was correct, but suits over executive conferences would make judges ``speech editors for communications issued by executive officials and event planners for meetings they hold.'' Officials must obey the Constitution even without judicial review.

\textbf{Scalia, with Thomas (concurring in the judgment).} The doctrine oscillates between ``Wallet Injury'' and ``Psychic Injury.'' Either psychic injury suffices and \textit{Flast} covers executive spending, or it does not and \textit{Flast} must go. The plurality's ``utterly meaningless distinctions'' amount to ``beating \textit{Flast} to a pulp'' while leaving it alive.

\textbf{Souter (dissenting).} Respondents sought only to apply \textit{Flast}. Under Madison's \textit{Memorial and Remonstrance}, the injury of having even ``three pence'' extracted for religion is the same whichever branch spends it; the distinction cannot be squared with \textit{Bowen}, and Article~III recognizes intangible harms.

\section{Examples: Future Applications}

\begin{enumerate}
\item \textbf{Same-Side: Agency Workshops.} An agency uses its general salaries account for workshops encouraging Christian charities to seek contracts. No specific enactment, so no taxpayer standing.
\item \textbf{Same-Side: Pentagon Booklets.} The Defense Department prints devotional booklets from general funds. \textit{Hein} bars taxpayer suit, though a proselytized visitor might sue.
\item \textbf{Opposite-Side: Specific Appropriation.} Congress appropriates \$50 million and directs grants to ``houses of worship and sectarian ministries.'' This satisfies \textit{Flast} and \textit{Bowen}.
\item \textbf{Opposite-Side: Disappointed Competitor.} A secular provider denied a grant because reviewers favor religious applicants has ordinary competitive-injury standing.
\item \textbf{Fence-Sitter: Committee-Report Earmark.} A \$200 million grant statute is silent on religion, but a committee report expects \$40 million for prison ministries. \textit{Bowen} suggests a specific program; Footnote~7 says earmarks lack legal force.
\end{enumerate}

\section{Critique}

\textbf{Originalist critique.} With Scalia, critics argue that Article~III's ``Cases'' and ``Controversies'' exclude generalized spending grievances and that the plurality's mandate line, unconnected to injury or redressability, was ``stealth overruling'' (Re) that left courts in a doctrine Judge Easterbrook called ``arbitrary'' and ``illogical.''

\textbf{Progressive critique.} Ira Lupu, Robert Tuttle, and others argue that because modern government runs on lump-sum appropriations and executive discretion, \textit{Hein} creates a large loophole in Establishment Clause enforcement and lets Congress shield religious programs by leaving allocation to the Executive. \textit{Winn} (2011) shrank \textit{Flast} further.

\section{Key Quotations}

\begin{quote}
``It is a necessary concomitant of the doctrine of \textit{stare decisis} that a precedent is not always expanded to the limit of its logic. That was the approach that then-Justice Rehnquist took in his opinion for the Court in \textit{Valley Forge}, and it is the approach we take here. We do not extend \textit{Flast}, but we also do not overrule it. We leave \textit{Flast} as we found it.''
\end{quote}
\noindent --- Justice Alito (plurality opinion), 551 U.S. at 615.

\begin{quote}
``In light of the size of the federal budget, it is a complete fiction to argue that an unconstitutional federal expenditure causes an individual federal taxpayer any measurable economic harm. And if every federal taxpayer could sue to challenge any Government expenditure, the federal courts would cease to function as courts of law and would be cast in the role of general complaint bureaus.''
\end{quote}
\noindent --- Justice Alito (plurality opinion), 551 U.S. at 593.

\begin{quote}
``If this Court is to decide cases by rule of law rather than show of hands, we must surrender to logic and choose sides: Either \textit{Flast v. Cohen}, 392 U. S. 83 (1968), should be applied to (at a minimum) all challenges to the governmental expenditure of general tax revenues in a manner alleged to violate a constitutional provision specifically limiting the taxing and spending power, or \textit{Flast} should be repudiated.''
\end{quote}
\noindent --- Justice Scalia (concurring in the judgment), 551 U.S. at 618.

\end{document}
